\renewcommand{\thechapter}{\Roman{chapter}}
\chapter{PROPOSED METHODOLOGY}
\renewcommand{\thechapter}{\arabic{chapter}}

% --- SECTION 3.1 ---
\section{Data Gathering}

% --- SUBSECTION 3.1.1 ---
\subsection{Dataset Source Selection}
	\begin{indentedblock}{4em}{0pt}
		Publicly available datasets for CAD file designs can be found on the internet. One prominent example is the ABC dataset, which consists of over one million unannotated STEP files derived from Onshape \parencite{abc_dataset_koch}. Another is the Fusion 360 Gallery, a dataset of free user-submitted designs with design timeline \parencite{fusion_dataset_willis}, and Text2CAD, which hosts around 170,000 CAD models with 660,000 natural language annotations ranging from short descriptions to detailed explanations \parencite{text2cad_mohammad}. Other sources include free 3D/CAD design websites such as GrabCAD and TraceParts, which contain various CAD objects not found in compiled datasets. A comparative summary for the sources listed above is contained in \autoref{tab:dataset_comparison_table} below.
		
		\skiplines{1}
		
		\begin{table}[!ht]
			\centering
			\begin{tblr}{
					cells={valign=m,halign=c},
					colspec={Q[1.8cm]Q[2cm]Q[4.5cm]Q[3cm]Q[3cm]},
					hlines,
					vlines
				}
				Source & Type & Summary & Advantages & Disadvantages \\
				TraceParts & {Online \\ Website} & {A library of standard engineering parts (bolts, gears, bearings) verified by manufacturers.} & {High-fidelity geometry, standard-compliant naming conventions.} & {Limited to standard parts; requires scraping.} \\
				GrabCAD & {Online \\ Website} & {Community repository with 6+ million user-generated CAD files.} & {High diversity of complex assemblies and creative designs.} & {Inconsistent quality; requires filtering for non-manifold geometry.} \\ 
				Text2CAD & {Compiled \\ Dataset} & {A collection of CAD files of different engineering parts with standardized annotation.} & {Ready-to-use text-geometry pairs; strictly organized.} & {Limited geometric complexity compared to industry files.} \\
				Fusion 360 & {Compiled \\ Dataset} & {Gallery of designs exported from Autodesk Fusion 360, focusing on assembly structures.} & {Contains reconstruction history (timelines).} & {Proprietary format issues; inconsistent export quality.} \\
				ABC & {Compiled \\ Dataset} & {Large-scale collection of 1M+ unlabelled B-Rep files from Onshape.} & {Massive scale useful for pre-training; robust topology.} & {Lack of semantic labels; contains many 'garbage' or tutorial files.} \\
				
			\end{tblr}
			\caption{Comparison between various dataset sources.}
			\label{tab:dataset_comparison_table}
		\end{table} 
		
		Based on the analysis detailed in \autoref{tab:dataset_comparison_table}, this paper will utilize a combination of CAD designs contained in Text2CAD as well as other online sources such as the ABC dataset and TraceParts. This combination ensures a balance between labeled semantic data (Text2CAD) and high-fidelity geometric variety (TraceParts). To help minimize discrepancies between annotations of several different datasets, a standard for annotation format mimicking the existing annotations in the Text2CAD dataset will be used. A summary of the procedure can be seen in \nameref{fig:filtering_flowchart} 
	\end{indentedblock}

\subsection{Dataset Preprocessing}
	\subsubsection{Dataset Filtering}
		\begin{indentedblock}{4em}{0pt}
			Before training, the dataset must be analyzed and filtered to remove low-quality CAD designs that could negatively affect the model's output. The analysis of the dataset is managed using Python with the \texttt{step-utils} and \textbf{Pandas} modules. A critical preprocessing step is \textbf{Manifold Validation}. As STEP files may contain disjointed surfaces or unstitched faces, any file failing a "water-tight" topological check (i.e., containing open shells) is discarded to prevent such outliers from polluting the model. Additional steps such as ensuring a balanced dataset (similar number of instances per class) to ensure statistical robustness and prevent model overfitting. More specifically, the class distribution is analyzed to identify and prune 'long-tail' categories, i.e. classes containing fewer than 100 unique geometric instances. Furthermore, deep learning models often resort to 'memorization' rather than generalization to handle classes with insufficient sample sizes \parencite{deep_learning_goodfellow}. As demonstrated by \parencite{generalization_zhang}, deep neural networks possess the capacity to fit arbitrary noise, meaning that in the absence of large-scale data constraints, the model tends to memorize specific training instances. This leads to 'shortcut learning,' where the model identifies objects via spurious noise patterns or texture statistics rather than learning the underlying topological features that define the class \parencites{shortcut_learning_geirhos}{memorization_arpit}. Research also shows that machine learning models struggle with identifying classes with low occurrences/high variability where the class pattern breaks \parencite{memorization_effect_you}. By enforcing a minimum threshold of $N=100$, we ensure that every semantic category (e.g., \texttt{GEAR}, \texttt{BRACKET}) possesses enough geometric variance to allow the Transformer to learn an abstract, generalized representation of the object type. The final dataset distribution is detailed in \autoref{sec:model_design}.
			
			\skiplines{1} 
			
		\end{indentedblock}

	\subsubsection{Annotation Format}
	\begin{indentedblock}{4em}{0pt}
		Since raw STEP files from sources like ABC do not contain semantic labels, a semi-automated annotation pipeline is employed. For TraceParts data, part names (e.g., "Hex Head Screw ISO 4017") are parsed and mapped to a simplified taxonomy of classes (e.g., \texttt{FASTENER}, \texttt{GEAR}, \texttt{BRACKET}). For Text2CAD, existing JSON-based annotations are retained. The final target label $Y$ for any given input $X$ is a categorical token representing the functional class of the 3D object.
	\end{indentedblock}

% --- SECTION 3.2 ---
\section{Model Design}
\label{sec:model_design}

This research proposes a novel \textit{Hierarchical Dual-Stream Tokenizer} designed specifically to address the non-Euclidean nature of B-Rep (Boundary Representation) data. Unlike standard Natural Language Processing (NLP) tokenizers that treat input as a 1D string, the proposed method treats the STEP file as a topological graph that must be normalized, compressed, and linearized before processing.

% --- SUBSECTION 3.2.1 ---
\subsection{Tokenizer Architecture}
	\begin{indentedblock}{4em}{0pt}
		The tokenization pipeline consists of four distinct stages designed to ensure geometric invariance and computational efficiency.
	\end{indentedblock}
	\subsubsection{Stage 1: Semantic Macro-Compression}
		\begin{indentedblock}{4em}{0pt}
			Raw STEP files suffer from extreme verbosity, where simple features (e.g., a hole) are described by dozens of low-level surfaces. To mitigate this, a graph contraction algorithm is applied. A rule-based parser scans the B-Rep graph for known topological sub-structures. For instance, a \texttt{CYLINDRICAL\_SURFACE} bounded by two \texttt{CIRCLE} edges is collapsed into a single macro-node \texttt{<CYLINDER\_PRIMITIVE>}. This reduces the sequence length by approximately 80\% while retaining semantic meaning.
		\end{indentedblock}
	
	\subsubsection{Stage 2: Canonical Geometric Normalization}
	\begin{indentedblock}{4em}{0pt}
		Standard STEP coordinates are sensitive to global positioning; a bolt at $x=0$ has different coordinate tokens than a bolt at $x=100$. To achieve translational and rotational invariance, the B-Rep graph undergoes \textbf{Canonical Alignment}:
		\begin{enumerate}
			\item \textbf{Centering:} The centroid of the solid is calculated, and all control points $P_i$ are translated such that the centroid lies at $(0,0,0)$.
			\item \textbf{Alignment:} Principal Component Analysis (PCA) is performed on the vertex cloud. The entire geometry is rotated via a rotation matrix $R$ such that the first principal component aligns with the global Z-axis.
		\end{enumerate}
	\end{indentedblock}
	
	\subsubsection{Stage 3: Topological Linearization (Canonical DFS)}
		\begin{indentedblock}{4em}{0pt}
			Since Transformers utilize positional encoding to interpret sequence order, the serialization of the B-Rep graph must be deterministic. A standard Depth-First Search (DFS) is insufficient because it is sensitive to the arbitrary indexing of nodes in the STEP file (e.g., visiting Face \#10 before Face \#20 purely based on ID).
			
			To resolve this, we implement a \textbf{Canonical DFS}. This algorithm introduces a strict ordering constraint at every branch of the traversal. For any parent node $N_p$ with a set of child nodes $C = \{c_1, c_2, ..., c_k\}$, the children are sorted based on a geometric invariant function $f(c)$ before traversal:
			
			\begin{equation}
				\text{Order}(C) = \text{sort}_{\downarrow}\left( \{ f(c) \mid c \in C \} \right)
			\end{equation}
			
			In this research, the invariant function $f(c)$ is defined as the \textbf{Surface Area} for topological faces and \textbf{Curve Length} for edges. This ensures that a specific 3D geometry yields an identical token sequence regardless of the permutation of entity IDs in the raw file, effectively removing graph isomorphism ambiguity from the input data.
		\end{indentedblock}
	
	\subsubsection{Stage 4: Complex-Valued Dual-Stream Encoding}
		\begin{indentedblock}{4em}{0pt}
			To preserve high-fidelity precision (e.g., NURBS weights) without bloating the vocabulary, a Dual-Stream architecture is used. The tokenizer emits two parallel streams for every step $t$:
			\begin{itemize}
				\item \textbf{Stream A (Semantic Tokens):} Discrete integers representing the type of entity (e.g., \texttt{<FACE>}, \texttt{<NURBS\_CURVE>}).
				\item \textbf{Stream B (Geometric Embeddings):} A 128-dimensional **Complex-Valued Vector** encoding the shape parameters.
			\end{itemize}
		\end{indentedblock}
	
	\begin{indentedblock}{4em}{0pt}
		The complex vector $v \in \mathbb{C}^{64}$ encodes the geometry as $z = r + i\theta$, where the real part $r$ represents scale-invariant magnitudes (e.g., radius, edge length) and the imaginary part $\theta$ encodes relative orientation using Rotary Positional Embeddings (RoPE).
	\end{indentedblock}

% --- SUBSECTION 3.2.2 ---
\subsection{Model Parameters}
	\begin{indentedblock}{4em}{0pt}
		The core model is a Transformer Encoder based on the BERT architecture, modified to accept the Dual-Stream input.
	
		\begin{itemize}
			\item \textbf{Embedding Dimension ($d_{model}$):} 512. The discrete tokens are embedded into $\mathbb{R}^{512}$, and the 128-dim complex geometric vectors are projected via a Multi-Layer Perceptron (MLP) to $\mathbb{R}^{512}$ before summation.
			\item \textbf{Attention Heads:} 8 Parallel heads.
			\item \textbf{Layers:} 12 Encoder blocks.
			\item \textbf{Context Window:} 2048 tokens (sufficient due to the Macro-Compression stage).
			\item \textbf{Feed-Forward Network:} GELU activation with a dimension of 2048.
		\end{itemize}
		
		The selection of these hyperparameters is governed by the specific trade-offs inherent to B-Rep graph learning:
		
		While LLMs often utilize dimensions upwards of 4096, CAD syntax has a significantly smaller vocabulary size ($<1000$ unique structural tokens) compared to natural language. A dimension of $d_{model}=512$ was chosen to prevent overfitting on the sparse vocabulary while providing sufficient manifold capacity to represent the fused Dual-Stream input. Specifically, it ensures that the projection of the 128-dimensional complex geometric vector does not dominate the semantic signal of the discrete token embedding during the summation process.
		
		The choice of 12 layers aligns with the hierarchical nature of boundary representation. Lower layers ($L_{1-4}$) are tasked with learning local topological syntax (e.g., edge continuity), intermediate layers ($L_{5-8}$) aggregate these into feature-level representations (e.g., pockets, bosses), and upper layers ($L_{9-12}$) extract high-level semantic class information. The 8 parallel attention heads allow the model to disentangle the multi-modal input: specific heads can learn to attend solely to the "Real" component of the input (spatial magnitude) while others attend to the "Imaginary" component (orientation) and topological connectivity, mimicking a multi-view perception system.
		
		Standard BERT implementations utilize a 512-token window. However, B-Rep graphs—even after the proposed Macro-Compression—exhibit long-range dependencies where a face defined at token $T_1$ may reference a surface defined at token $T_{1000}$. A window of 2048 tokens was empirically selected to encompass the complete topological graph of 95\% of the dataset samples without necessitating destructive truncation, striking a balance between memory constraints and geometric completeness.

	\end{indentedblock}
	
	\subsubsection{Geometry-Infused Attention Mechanism}
		\begin{indentedblock}{4em}{0pt}
			Unlike standard NLP Transformers where the attention score is derived solely from semantic word embeddings, the proposed architecture employs a \textbf{Geometry-Infused Attention} mechanism. This mechanism is responsible for synthesizing the Dual-Stream output from the tokenizer into a unified representation.
			
			The input to the attention layer for the $i$-th token is defined as the fusion of the semantic and geometric streams:
			\begin{equation}
				X_i = \text{Embed}(T_i) + \text{MLP}(V_i)
			\end{equation}
			Where $T_i$ is the discrete semantic token (e.g., \texttt{<FACE>}) and $V_i$ is the 128-dimensional complex geometric vector. The Multi-Layer Perceptron (MLP) projects the continuous geometric manifold into the model dimension $d_{model}$.
			
			\paragraph{Rotary Positional Integration}
			To strictly enforce the rotational invariance defined in the methodology, the model utilizes a modified version of Rotary Positional Embeddings (RoPE). The "Phase" component ($\theta$) of the complex geometric input $V_i$ is directly mapped to the rotation of the Query ($Q$) and Key ($K$) vectors in the attention calculation:
			\begin{equation}
				\text{Attention}(Q, K, V) = \text{softmax}\left(\frac{(R_{\theta_i}Q)(R_{\theta_j}K)^T}{\sqrt{d_k}}\right)V
			\end{equation}
			Here, $R_{\theta}$ represents the rotation matrix derived from the tokenizer's complex input. This ensures that the attention score between two CAD entities depends not only on their semantic type but also on their \textbf{relative geometric orientation}. For example, a \texttt{SCREW} entity will attend strongly to a \texttt{HOLE} entity only if their orientation vectors (Normal vectors) are aligned (i.e., the relative rotation is near zero), thereby embedding physical assembly logic directly into the attention weights.
		\end{indentedblock}
	
	[Image of multi-head attention architecture diagram]
	
	\begin{indentedblock}{4em}{0pt}
		The model output is passed through a global average pooling layer followed by a softmax classifier corresponding to the number of classes defined in the dataset preprocessing stage.
		
	\end{indentedblock}

% --- SECTION 3.3 ---
\section{Training and Backtesting}

	\subsection{Training Setup}
		\begin{indentedblock}{4em}{0pt}
			The classification objective is optimized using \textbf{Categorical Cross-Entropy Loss}. This function penalizes the divergence between the predicted probability distribution $\hat{y}$ and the ground truth one-hot vector $y$. For a dataset of $N$ samples and $C$ classes, the loss $\mathcal{L}$ is defined as:
			
			\begin{equation}
				\mathcal{L} = -\frac{1}{N} \sum_{i=1}^{N} \sum_{c=1}^{C} y_{i,c} \log(\hat{y}_{i,c})
			\end{equation}
			
			Where $\hat{y}_{i,c}$ is the softmax probability output of the model for class $c$. The logarithmic nature of the loss function imposes a heavy penalty on confident but incorrect predictions, which is critical for distinguishing between visually similar CAD classes (e.g., distinguishing an ISO-4014 Bolt from an ISO-4017 Bolt based on subtle shaft length variations).
		\end{indentedblock}

\subsection{Evaluation Metrics}
	\begin{indentedblock}{4em}{0pt}
		To validate the effectiveness of the proposed tokenizer, the model is evaluated on two fronts:
		\begin{enumerate}
			\item \textbf{Standard Classification Metrics:} Accuracy, Precision, Recall, and F1-Score on a held-out test set (20\% of the data).
			\item \textbf{Rotational Robustness Test:} A specialized test set is generated by applying random random $SO(3)$ rotations to the test data. The model's accuracy on this rotated set is compared to the baseline to verify the efficacy of the Canonical Alignment and Complex-Valued Embeddings.
		\end{enumerate}
	\end{indentedblock}